\chapter{Durchführung} \label{ch:Durchfuehrung}

Auch der Folgende Text ist am Praktikumsskript von Jens Wagner orientiert. Es gelten diesleben Hinweise wie in \cch{Einleitung}.\cite{skriptPAP21,Duden}

\section{Messverfahren}

Die Messungen wurden mit einem Gitterspektrometer des Typs \emph{Ocean Optics USB4000} durchgeführt. Für jede Lichtquelle wurde die Integrationszeit so gewählt, dass die relevanten Spektralbereiche gut aufgelöst und nicht gesättigt waren. Zusätzlich wurde die Mittelungsfunktion der Software \emph{OceanView} genutzt, um das Rauschen zu reduzieren. Dabei durfte die Gesamtmesszeit pro Aufnahme vier Sekunden nicht überschreiten.

Im ersten Versuchsteil wurde das Sonnenspektrum aufgenommen. Dazu wurde der Lichteinfall einmal durch das geöffnete Fenster und einmal durch das geschlossene Fenster gemessen. Ein direktes Ausrichten des Spektrometers auf die Sonne war nicht möglich, sodass das diffuse Himmelslicht verwendet wurde.

Im zweiten Versuchsteil wurden verschiedene künstliche Lichtquellen untersucht, darunter Glühlampen, Energiesparlampen, LEDs unterschiedlicher Farben sowie ein Laser. Die Lichtquellen wurden jeweils vor den Eingang des Spektrometers positioniert. Abstand und Blendenöffnung wurden so angepasst, dass ein möglichst großer, aber nicht gesättigter Intensitätsbereich erfasst wurde.

Im dritten Versuchsteil wurde das Emissionsspektrum einer Natriumdampflampe aufgenommen. Um sowohl starke als auch schwache Linien sichtbar zu machen, wurden mehrere Messungen mit unterschiedlichen Integrationszeiten und Blendenstellungen durchgeführt. Für die Analyse wurden jeweils die Aufnahmen verwendet, in denen die betreffenden Linien am deutlichsten erkennbar waren.

\section{Versuchsaufbau}

Der Versuchsaufbau besteht aus dem Gitterspektrometer, einem optischen Eingang mit verstellbarer Blende sowie einem PC zur Datenerfassung. Die Lichtquellen wurden so positioniert, dass ihr Licht direkt in den Eingang des Spektrometers gelangte. Bei der Messung des Sonnenlichts wurde das Spektrometer auf den Himmel ausgerichtet, wobei der Einfluss des Fensterglases durch Vergleichsmessungen bestimmt wurde.

Für die Natriumdampflampe wurde ein sicherer Abstand eingehalten, um eine Überbelichtung des Detektors zu vermeiden. Durch Variation der Integrationszeit konnten sowohl die intensiven Linien der D-Dubletts als auch die schwächeren Linien der Nebenserien erfasst werden. Die Rohdaten wurden anschließend gespeichert und für die Auswertung weiterverarbeitet.
